\ifdefined\gkver\else\def\gkver{student}\fi
\documentclass[\gkver]{gaokaozhenti}
\begin{document}

\chapter{2016年北京理综}
%% paperType: 理综物理部分

\begin{enumerate}
\item
%% number: 13
%% paperName: 2016年北京理综第 13 题 6 分
%% typeId: 1
%% type: 单选题
%% chapter: 原子和原子核
%% point: 玻尔模型
%% method: 无
%% score: 6
%% degree: 850
%% duplicateId: 0
%% body:
处于 $n=3$ 能级的大量氢原子，向低能级跃迁时，辐射光的频率有\xzanswer{C}
\fourchoices[answer=C]
{1 种}
{2 种}
{3 种}
{4 种}
%% answer: C
%% memo:
\memoanswer{
处于 $n=3$ 能级的大量氢原子，向低能级跃迁时，会出现 $n=3$ 向 $n=2$、$n=1$ 直接跃迁的两种情形，辐射出两种频率的光子，还有 $n=2$ 向 $n=1$ 跃迁时再辐射出一种光子，所以会出现 3 种频率的光子，故 C 正确。
}

\item
%% number: 14
%% paperName: 2016年北京理综第 14 题 6 分
%% typeId: 1
%% type: 单选题
%% chapter: 光学
%% point: 光的电磁说
%% method: 无
%% score: 6
%% degree: 750
%% duplicateId: 0
%% body:
下列说法正确的是\xzanswer{A}
\fourchoices[answer=A]
{电磁波在真空中以光速 $c$ 传播}
{在空气中传播的声波是横波}
{声波只能在空气中传播}
{光需要介质才能传播}
%% answer: A
%% memo:
\memoanswer{
电磁波在真空中以光速 $c$ 传播，故 A 正确；在空气中传播的声波是纵波，故 B 错误；声波在气体、液体、固体中都可以传播，故 C 错误；光是电磁波，传播不需要介质，故 D 错误。
}

\item
%% number: 15
%% paperName: 2016年北京理综第 15 题 6 分
%% typeId: 1
%% type: 单选题
%% chapter: 机械振动
%% point: 振动图象
%% method: 无
%% score: 6
%% degree: 750
%% duplicateId: 0
%% body:
如图所示，弹簧振子在 M、N 之间做简谐运动。以平衡位置 $O$ 为原点，建立 $Ox$ 轴，向右为 $x$ 轴正方向。若振子位于 N 点时开始计时，则其振动图像为\xzanswer{A}

\onepicture[w=6cm]{\includesvg{figs/15}}

\fourchoices[answer=A, ispicture=true, h=2.4cm]
{figs/15a.png}
{figs/15b.png}
{figs/15c.png}
{figs/15d.png}
%% answer: A
%% memo:
\memoanswer{
振动图像以 $O$ 为平衡位置、向右为 $x$ 轴正方向，从振子位于 N 点时开始计时，即 0 时刻振子的位移为正向最大值，故 A 正确。
}

\item
%% number: 16
%% paperName: 2016年北京理综第 16 题 6 分
%% typeId: 1
%% type: 单选题
%% chapter: 电磁感应
%% point: 法拉第电磁感应定律
%% method: 无
%% score: 6
%% degree: 650
%% duplicateId: 0
%% body:
如图所示，匀强磁场中有两个导体圆环 a、b，磁场方向与圆环所在平面垂直。磁感应强度 $B$ 随时间均匀增大。两圆环半径之比为 $2:1$，圆环中产生的感应电动势分别为 $E_{a}$ 和 $E_{b}$。不考虑两圆环间的相互影响。下列说法正确的是\xzanswer{B}

\onepicture[w=6cm]{figs/16.png}

\fourchoices[answer=B]
{$E_{a}:E_{b}=4:1$，感应电流均沿逆时针方向}
{$E_{a}:E_{b}=4:1$，感应电流均沿顺时针方向}
{$E_{a}:E_{b}=2:1$，感应电流均沿逆时针方向}
{$E_{a}:E_{b}=2:1$，感应电流均沿顺时针方向}
%% answer: B
%% memo:
\memoanswer{
由法拉第电磁感应定律得 $E=\frac{\Delta\phi}{\Delta t}=\frac{\Delta B\cdot S}{\Delta t}$，所以圆环的感应电动势之比为 $\frac{E_{a}}{E_{b}}=\frac{S_{a}}{S_{b}}$，又圆环的面积 $S=\pi r^{2}$，所以 $\frac{E_{a}}{E_{b}}=\frac{S_{a}}{S_{b}}=\frac{r_{a}^{2}}{r_{b}^{2}}=\frac{4}{1}$；因为磁场方向垂直纸面向外且磁感应强度增大，由楞次定律可以判断圆环内的感应电流沿顺时针方向，故 B 正确。
}

\item
%% number: 17
%% paperName: 2016年北京理综第 17 题 6 分
%% typeId: 1
%% type: 单选题
%% chapter: 磁场
%% point: 磁感线
%% method: 无
%% score: 6
%% degree: 550
%% duplicateId: 0
%% body:
中国宋代科学家沈括在《梦溪笔谈》中最早记载了地磁偏角：“以磁石磨针锋，则能指南，然常微偏东，不全南也。”进一步研究表明，地球周围地磁场的磁感线分布示意如图。结合上述材料，下列说法不正确的是\xzanswer{C}

\onepicture[w=6cm]{figs/17.png}

\fourchoices[answer=C]
{地理南、北极与地磁场的南、北极不重合}
{地球内部也存在磁场，地磁南极在地理北极附近}
{地球表面任意位置的地磁场方向都与地面平行}
{地磁场对射向地球赤道的带电宇宙射线粒子有力的作用}
%% answer: C
%% memo:
\memoanswer{
地磁偏角的存在说明地理南、北极与地磁场的南、北极不重合，故 A 正确；由地球周围地磁场的磁感线分布示意图及磁感线是闭合曲线，可知地球内部也存在磁场，地磁南极在地理北极附近，故 B 正确；由图可知，在地球的两极附近，磁感线和地面是倾斜相交的，不与地面平行，故 C 错误；带电粒子速度与磁感线不平行时会受到洛伦兹力的作用，故 D 正确。
}

\item
%% number: 18
%% paperName: 2016年北京理综第 18 题 6 分
%% typeId: 1
%% type: 单选题
%% chapter: 曲线运动
%% point: 人造地球卫星
%% method: 无
%% score: 6
%% degree: 450
%% duplicateId: 0
%% body:
如图所示，一颗人造卫星原来在椭圆轨道 1 绕地球 E 运行，在 P 点变轨后进入轨道 2 做匀速圆周运动。下列说法正确的是\xzanswer{B}

\onepicture[w=5.5cm]{\includesvg{figs/18}}

\fourchoices[answer=B]
{不论在轨道 1 还是轨道 2 运行，卫星在 P 点的速度都相同}
{不论在轨道 1 还是轨道 2 运行，卫星在 P 点的加速度都相同}
{卫星在轨道 1 的任何位置都具有相同加速度}
{卫星在轨道 2 的任何位置都具有相同动量}
%% answer: B
%% memo:
\memoanswer{
卫星在 P 点变轨，能够进入轨道 2，说明卫星相对轨道 1 在做离心运动，所以卫星在轨道 2 上 P 点的速度大于在轨道 1 上 P 点的速度，故 A 错误；无论在哪个轨道上，卫星在 P 点受到的万有引力相同，由 $\frac{GMm}{r^{2}}=ma$ 可知，其在 P 点的加速度相同，故 B 正确；万有引力的方向始终指向地心，所以加速度的方向也始终指向地心，时刻在变化，故 C 错误；卫星在轨道 2 上的速度方向沿着该点的切向，时刻在发生变化，所以动量的方向也是时刻变化的，故 D 错误。
}

\item
%% number: 19
%% paperName: 2016年北京理综第 19 题 6 分
%% typeId: 1
%% type: 单选题
%% chapter: 电路
%% point: 测电动势与内阻
%% method: 无
%% score: 6
%% degree: 450
%% duplicateId: 0
%% body:
某兴趣小组探究用不同方法测定干电池的电动势和内阻，他们提出的实验方案中有如下四种器材组合。为使实验结果尽可能准确，最不可取的一组器材是\xzanswer{D}
\fourchoices[answer=D]
{一个安培表、一个伏特表和一个滑动变阻器}
{一个伏特表和多个定值电阻}
{一个安培表和一个电阻箱}
{两个安培表和一个滑动变阻器}
%% answer: D
%% memo:
\memoanswer{
一个安培表、一个伏特表和一个滑动变阻器，利用公式 $E=U+Ir$，测量两组数据，可得电池的电动势和内阻，故 A 可取；一个伏特表和多个定值电阻，利用公式 $E=U+\frac{U}{R}r$，也可以得到电动势和内阻，故 B 可取；一个安培表和一个电阻箱，利用公式 $E=IR+Ir$，同理可得电动势和内阻，故 C 可取；两个安培表和一个滑动变阻器，只能测量电流，无法测量电压和电阻，所以无法得到电动势和内阻，故 D 不可取。
}

\item
%% number: 20
%% paperName: 2016年北京理综第 20 题 6 分
%% typeId: 1
%% type: 单选题
%% chapter: 气体性质
%% point: 分子动理论
%% method: 无
%% score: 6
%% degree: 550
%% duplicateId: 0
%% body:
雾霾天气是对大气中各种悬浮颗粒物含量超标的笼统表述，是特定气候条件与人类活动相互作用的结果。雾霾中，各种悬浮颗粒物形状不规则，但可视为密度相同、直径不同的球体，并用 PM10、PM2.5 分别表示球体直径小于或等于 $10\Uum$、$2.5\Uum$ 的颗粒物（PM 是颗粒物的英文缩写）。

某科研机构对北京地区的检测结果表明，在静稳的雾霾天气中，近地面高度百米的范围内，PM10 的浓度随高度的增加略有减小，大于 PM10 的大悬浮颗粒物的浓度随高度的增加明显减小，且两种浓度分布基本不随时间变化。

据此材料，以下叙述正确的是\xzanswer{C}
\fourchoices[answer=C]
{PM10 表示直径小于或等于 $1.0\times10^{-6}\Um$ 的悬浮颗粒物}
{PM10 受到的空气分子作用力的合力始终大于其所受到的重力}
{PM10 和大悬浮颗粒物都在做布朗运动}
{PM2.5 的浓度随高度的增加逐渐增大}
%% answer: C
%% memo:
\memoanswer{
PM$_{10}$ 表示直径小于或等于 $10\Uum=1.0\times10^{-5}\Um$，故 A 错误；PM$_{10}$ 颗粒悬浮在空气中，其运动状态不能确定，故其受力情况不能确定，故 B 错误；材料中指出 PM$_{10}$ 和大悬浮颗粒物，两种浓度分布基本不随时间变化，说明两种颗粒在做无规则运动，故两种颗粒都在做布朗运动，故 C 正确；材料中并没有说明 PM$_{2.5}$ 浓度随高度的增加逐渐增大，故 D 错误。
}

\item
%% number: 21
%% paperName: 2016年北京理综第 21 题 12 分
%% typeId: 4
%% type: 实验题
%% chapter: 功和能
%% point: DIS研究机械能守恒定律
%% method: 无
%% score: 12
%% degree: 450
%% duplicateId: 0
%% body:
\begin{enumerate}
	\item
	热敏电阻常用于温度控制或过热保护装置中。如图为某种热敏电阻和金属热电阻的阻值 $R$ 随温度 $t$ 变化的示意图。由图可知，这种热敏电阻在温度上升时导电能力 \tkanswer[1.5cm]{增强}（选填“增强”或“减弱”）；相对金属热电阻而言，热敏电阻对温度变化的响应更 \tkanswer[1.5cm]{敏感}（选填“敏感”或“不敏感”）。
	\onepicture[h=4cm, align=right]{\includesvg{figs/21}}

	\item
	利用如图 \ref{2016北京理综21a} 所示装置做“验证机械能守恒定律”实验。
	\begin{steps}
		\item
		为验证机械能是否守恒，需要比较重物下落过程中任意两点间的 \tkanswer[1cm]{A}。

		A．动能变化量与势能变化量

		B．速度变化量与势能变化量

		C．速度变化量与高度变化量

		\item
		除带夹子的重物、纸带、铁架台（含铁夹）、电磁打点计时器、导线及开关外，在下列器材中，还必须使用的两种器材是 \tkanswer[1cm]{AB}。

		A．交流电源　　　　B．刻度尺　　　　C．天平（含砝码）

		\item
		实验中，先接通电源，再释放重物，得到如图 \ref{2016北京理综21b} 所示的一条纸带。在纸带上选取三个连续打出的点 A、B、C，测得它们到起始点 O 的距离分别为 $h_{A}$、$h_{B}$、$h_{C}$。

		已知当地重力加速度为 $g$，打点计时器打点的周期为 $T$。设重物的质量为 $m$。从打 O 点到打 B 点的过程中，重物的重力势能变化量 $\Delta E_{p}=$ \tkanswer[2cm]{$-mgh_{B}$}，动能变化量 $\Delta E_{k}=$ \tkanswer[3cm]{$\frac{1}{2}m\left(\frac{h_{C}-h_{A}}{2T}\right)^{2}$}。

		\item
		大多数学生的实验结果显示，重力势能的减少量大于动能的增加量，原因是 \tkanswer[1cm]{C}。

		A．利用公式 $v=gt$ 计算重物速度

		B．利用公式 $v=\sqrt{2gh}$ 计算重物速度

		C．存在空气阻力和摩擦阻力的影响

		D．没有采用多次实验取平均值的方法

		\item
		某同学想用下述方法研究机械能是否守恒：在纸带上选取多个计数点，测量它们到起始点 O 的距离 $h$，计算对应计数点的重物速度 $v$，描绘 $v^{2}-h$ 图像，并做如下判断：若图像是一条过原点的直线，则重物下落过程中机械能守恒。请你分析论证该同学的判断依据是否正确。
	\end{steps}
	\twopicture[labelA=2016北京理综21a, labelB=2016北京理综21b, numA=1, numB=2, align=right, wA=4.5cm, wB=8cm]
	{figs/21a.png}
	{figs/21b.png}
\end{enumerate}
%% answer: 
\jdanswer{
\begin{enumerate}
	\item
	增强，敏感

	\item
	①A；②AB；③$-mgh_{B}$，$\frac{1}{2}m\left(\frac{h_{C}-h_{A}}{2T}\right)^{2}$；④C；⑤该同学的判断依据不正确。在重物下落 $h$ 的过程中，若阻力 $f$ 恒定，由 $mgh-fh=\frac{1}{2}mv^{2}-0$ 得 $v^{2}=2\left(g-\frac{f}{m}\right)h$ 可知，$v^{2}-h$ 图像就是过原点的一条直线。要想通过 $v^{2}-h$ 图像的方法验证机械能是否守恒，还必须看图像的斜率是否接近 $2g$。

\end{enumerate}
}
%% memo:
\memoanswer{
（1）由图可知，热敏电阻的阻值随温度的升高而减小，电阻减小则其导电能力增强；由图可知，在相同的温度变化范围内，热敏电阻的阻值变化更大，即热敏电阻对温度的变化更敏感。

（2）①如果机械能守恒，则动能和势能相互转化，所以本实验需要比较重物下落过程中两点的动能变化量与势能变化量，故 A 正确。

②电磁打点计时器需要交流电源，数据处理时需要利用刻度尺测量纸带上的距离，故 AB 正确；本实验不需要测量重物的质量，所以不需要天平。

③从 O 点到 B 点，重力势能的变化为 $\Delta E_{p}=0-mgh_{B}=-mgh_{B}$，因起始点的速度为零，所以动能的变化量等于 B 点的动能，B 的速度为 $v=\frac{h_{C}-h_{A}}{2T}$，则动能变化量为 $\Delta E_{k}=\frac{1}{2}mv^{2}=\frac{m(h_{C}-h_{A})^{2}}{8T^{2}}$。

④由于实验中存在空气阻力和摩擦阻力，重物实际速度小于理论速度，所以动能的增加量小于重力势能的减少量，故 C 正确。

⑤该同学的判断依据不正确。在重物下落 $h$ 的过程中，若阻力 $f$ 恒定，由 $mgh-fh=\frac{1}{2}mv^{2}-0$ 得 $v^{2}=2\left(g-\frac{f}{m}\right)h$ 可知，$v^{2}-h$ 图像就是过原点的一条直线。要想通过 $v^{2}-h$ 图像的方法验证机械能是否守恒，还必须看图像的斜率是否接近 $2g$。
}

\item
%% number: 22
%% paperName: 2016年北京理综第 22 题 16 分
%% typeId: 6
%% type: 计算题
%% chapter: 磁场
%% point: 带电粒子在磁场中的运动
%% method: 无
%% score: 16
%% degree: 650
%% duplicateId: 0
%% body:
如图所示，质量为 $m$、电荷量为 $q$ 的带电粒子，以初速度 $v$ 沿垂直磁场方向射入磁感应强度为 $B$ 的匀强磁场，在磁场中做匀速圆周运动。不计带电粒子所受重力。
\begin{enumerate}
	\item
	求粒子做匀速圆周运动的半径 $R$ 和周期 $T$；
	\item
	为使该粒子做匀速直线运动，还需要同时存在一个与磁场方向垂直的匀强电场，求电场强度 $E$ 的大小。
\end{enumerate}
\onepicture[align=right, w=4cm]{figs/22.png}
%% answer: 
\jdanswer{
\begin{enumerate}
	\item
	$R=\frac{mv}{qB}$，$T=\frac{2\pi m}{qB}$

	\item
	$E=vB$

\end{enumerate}
}
%% memo:
\memoanswer{
（1）带电粒子在磁场中，由洛伦兹力提供向心力得
$f=qvB=m\frac{v^{2}}{R}$，
带电粒子做匀速圆周运动的半径为
$R=\frac{mv}{qB}$，
匀速圆周运动的周期
$T=\frac{2\pi R}{v}=\frac{2\pi m}{qB}$。

（2）粒子受电场力 $F=qE$，洛伦兹力 $f=qvB$。粒子做匀速直线运动，则 $qE=qvB$，场强 $E$ 的大小 $E=vB$。
}

\item
%% number: 23
%% paperName: 2016年北京理综第 23 题 18 分
%% typeId: 6
%% type: 计算题
%% chapter: 电场
%% point: 带电粒子的偏转
%% method: 无
%% score: 18
%% degree: 450
%% duplicateId: 0
%% body:
如图所示，电子由静止开始经加速电场加速后，沿平行于板面的方向射入偏转电场，并从另一侧射出。已知电子质量为 $m$，电荷量为 $e$，加速电场电压为 $U_{0}$。偏转电场可看做匀强电场，极板间电压为 $U$，极板长度为 $L$，板间距为 $d$。
\begin{enumerate}
	\item
	忽略电子所受重力，求电子射入偏转电场时的初速度 $v_{0}$ 和从电场射出时沿垂直板面方向的偏转距离 $\Delta y$；
	\item
	分析物理量的数量级，是解决物理问题的常用方法。在解决（1）问时忽略了电子所受重力，请利用下列数据分析说明其原因。已知 $U=2.0\times10^{2}\UV$，$d=4.0\times10^{-2}\Um$，$m=9.1\times10^{-31}\Ukg$，$e=1.6\times10^{-19}\UC$，$g=10\Umsq$。
	\item
	极板间既有静电场也有重力场。电势反映了静电场各点的能的性质，请写出电势 $\varphi$ 的定义式。类比电势的定义方法，在重力场中建立“重力势” $\varphi_{G}$ 的概念，并简要说明电势和“重力势”的共同特点。
\end{enumerate}
\onepicture[align=right, w=8cm]{figs/23.png}
%% answer: 
\jdanswer{
\begin{enumerate}
	\item
	$v_{0}=\sqrt{\frac{2eU_{0}}{m}}$，$\Delta y=\frac{UL^{2}}{4U_{0}d}$

	\item
	重力 $G=mg\sim10^{-29}\UN$，电场力 $F=\frac{eU}{d}\sim10^{-15}\UN$，由于 $F\gg G$，因此不需要考虑电子所受重力。

	\item
	$\varphi=\frac{E_{p}}{q}$，$\varphi_{G}=\frac{E_{G}}{m}$；电势 $\varphi$ 和重力势 $\varphi_{G}$ 都是反映场的能的性质的物理量，仅由场自身的因素决定。

\end{enumerate}
}
%% memo:
\memoanswer{
（1）由功和能的关系得
$eU_{0}=\frac{1}{2}mv_{0}^{2}$，
电子射入偏转电场的初速度为
$v_{0}=\sqrt{\frac{2eU_{0}}{m}}$，
在偏转电场中，电子的运动时间为
$\Delta t=\frac{L}{v_{0}}=L\sqrt{\frac{m}{2eU_{0}}}$，
偏转距离为
$\Delta y=\frac{1}{2}a(\Delta t)^{2}=\frac{UL^{2}}{4U_{0}d}$。

（2）考虑电子所受重力和电场力的数量级，有
重力 $G=mg\sim10^{-29}\UN$，
电场力 $F=\frac{eU}{d}\sim10^{-15}\UN$，
由于 $F\gg G$，因此不需要考虑电子所受重力。

（3）电场中某点电势 $\varphi$ 定义为电荷在该点的电势能 $E_{p}$ 与其电荷量 $q$ 的比值，即
$\varphi=\frac{E_{p}}{q}$，
由于重力做功与路径无关，可以类比静电场电势的定义，将重力场中物体在某点的重力势能 $E_{G}$ 与其质量 $m$ 的比值，叫作“重力势”，即
$\varphi_{G}=\frac{E_{G}}{m}$，
电势 $\varphi$ 和重力势 $\varphi_{G}$ 都是反映场的能的性质的物理量，仅由场自身的因素决定。
}

\item
%% number: 24
%% paperName: 2016年北京理综第 24 题 10 分
%% typeId: 6
%% type: 计算题
%% chapter: 运动和力
%% point: 动量定理
%% method: 无
%% score: 10
%% degree: 450
%% duplicateId: 0
%% body:
\begin{enumerate}
	\item
	动量定理可以表示为 $\Delta p=F\Delta t$，其中动量 $p$ 和力 $F$ 都是矢量。在运用动量定理处理二维问题时，可以在相互垂直的 $x$、$y$ 两个方向上分别研究。例如，质量为 $m$ 的小球斜射到木板上，入射的角度是 $\theta$，碰撞后弹出的角度也是 $\theta$，碰撞前后的速度大小都是 $v$，如图 \ref{2016北京理综24a} 所示。碰撞过程中忽略小球所受重力。

	A．分别求出碰撞前后 $x$、$y$ 方向小球的动量变化量 $\Delta p_{x}$、$\Delta p_{y}$；

	B．分析说明小球对木板的作用力的方向。
	\onepicture[h=4cm, align=right, num=1, showlabel=false, label=2016北京理综24a]{figs/24a.png}

	\item
	激光束可以看作是粒子流，其中的粒子以相同的动量沿光传播方向运动。激光照射到物体上，在发生反射、折射和吸收现象的同时，也会对物体产生作用。光镊效应就是一个实例，激光束可以像镊子一样抓住细胞等微小颗粒。一束激光经 S 点后被分成若干细光束，若不考虑光的反射和吸收，其中光束①和②穿过介质小球的光路如图 \ref{2016北京理综24b} 所示。图中 O 点是介质小球的球心，入射时光束①和②与 SO 夹角均为 $\theta$，出射时光束均与 SO 平行。请在下面两种情况下，分析说明两光束因折射对小球产生的合力的方向。

	A．光束①和②强度相同；

	B．光束①比②的强度大。
	\onepicture[h=3.5cm, align=right, num=2, showlabel=false, label=2016北京理综24b]{figs/24b.png}
\end{enumerate}
%% answer: 
\jdanswer{
\begin{enumerate}
	\item
	A．$\Delta p_{x}=0$，$\Delta p_{y}=2mv\cos\theta$，方向沿 $y$ 轴正方向。

	B．根据动量定理可知，木板对小球作用力的方向沿 $y$ 轴正方向；根据牛顿第三定律可知，小球对木板作用力的方向沿 $y$ 轴负方向。

	\item
	A．两光束对小球的合力的方向沿 SO 向左。

	B．两光束对小球的合力的方向指向左上方。

\end{enumerate}
}
%% memo:
\memoanswer{
（1）a．$x$ 方向：动量变化为
$\Delta p_{x}=mv\sin\theta-mv\sin\theta=0$，
$y$ 方向：动量变化为
$\Delta p_{y}=mv\cos\theta-(-mv\cos\theta)=2mv\cos\theta$，
方向沿 $y$ 轴正方向。

b．由动量定理可知，木板对小球作用力的方向沿 $y$ 轴正方向；由牛顿第三定律可知，小球对木板作用力的方向沿 $y$ 轴负方向。

（2）a．仅考虑光的折射，设 $\Delta t$ 时间内每束光穿过小球的粒子数为 $n$，每个粒子动量的大小为 $p$。这些粒子进入小球前的总动量为
$p_{1}=2np\cos\theta$，
从小球出射时的总动量为
$p_{2}=2np$，
$p_{1}$、$p_{2}$ 的方向均沿 SO 向右，由动量定理得
$F\Delta t=p_{2}-p_{1}=2np(1-\cos\theta)>0$，
可知小球对这些粒子的作用力 $F$ 的方向沿 SO 向右；由牛顿第三定律，两光束对小球的合力的方向沿 SO 向左。

b．建立如图所示的 $Oxy$ 直角坐标系。

\onepicture[w=8cm]{figs/24_an.png}

$x$ 方向：由（2）a 同理可知，两光束对小球的作用力沿 $x$ 轴负方向。

$y$ 方向：设 $\Delta t$ 时间内，光束①穿过小球的粒子数为 $n_{1}$，光束②穿过小球的粒子数为 $n_{2}$，$n_{1}>n_{2}$。这些粒子进入小球前的总动量为
$p_{1y}=(n_{1}-n_{2})p\sin\theta$，
从小球出射时的总动量为
$p_{2y}=0$，
由动量定理得
$F_{y}\Delta t=p_{2y}-p_{1y}=-(n_{1}-n_{2})p\sin\theta$，
可知小球对这些粒子的作用力 $F_{y}$ 的方向沿 $y$ 轴负方向，由牛顿第三定律，两光束对小球的作用力沿 $y$ 轴正方向。所以两光束对小球的合力的方向指向左上方。
}

\end{enumerate}

\end{document}
